\chapter{Fisher's Estimation Criteria and Their Regularity Boundary}
\label{sec:p15-fisher-estimation}

Fisher's 1922 \emph{On the Mathematical Foundations of Theoretical
Statistics} separates a population parameter from the statistic calculated
to estimate it. He organizes statistical work into specifying a population
law, estimating its parameters, and finding the sampling distributions of
the resulting statistics. Maximizing a function addresses only part of that
program. Whether the maximizer exists, identifies the target, retains the
sample's information, or has a useful sampling approximation requires
additional argument \cite[pp.~309--331]{Fisher1922Foundations}.

The core reconstruction below follows that original presentation, including
its binomial coordinate comparison and its admitted gap in the general
sufficiency argument. The exact finite results are separately attributed
handbook derivations. The general sufficiency criterion, the consistency
and asymptotic normality theorems for maximum likelihood with their
regularity hypotheses, and the Type III and location--scale calculations
are in \cref{sec:p15-w02b-p15-057}.

\section{Population, sample and three different problems}

The sampling experiment is the family \(P_\theta\) in
\cref{def:p15-statistical-experiment}. A \emph{statistic} is a measurable
function of the recorded data, not of the unknown parameter. An
\emph{estimand} is the target function of that parameter. A likelihood
holds the observation fixed and compares the laws' densities or masses
at that observation. A sampling distribution holds the parameter fixed
and varies the data. The posterior in \cref{def:p15-bayes-experiment}
requires a further parameter probability measure.

Fisher's specification problem concerns which mathematical population
law the observations should represent. His estimation problem asks which
sample functions estimate its parameters. His distribution problem asks
how those sample functions vary in repeated sampling. The claim that
independent measurements can be interpreted as a random sample does not
establish identical distributions, independence, a correct parametric
family or a representative observation mechanism. Those are substantive
model premises. Nor does a goodness-of-fit comparison prove the chosen
family true \cite[pp.~311--315]{Fisher1922Foundations}.

His \emph{consistency} criterion initially says that applying the
statistical calculation to the whole hypothetical population should
return its parameter. This is a population-functional requirement.
Modern sampling consistency additionally asks whether an estimator
sequence converges in probability to the target under each stated law.
The population identity alone supplies neither existence of required
moments nor that convergence.

Efficiency in the paper compares large-sample probable errors when the
statistics approach normal laws with the specified sample-size scaling.
Sufficiency asks whether the statistic retains all information relevant
to the parameter: after fixing its value, another statistic's conditional
distribution should no longer depend on the parameter. These are not
synonyms for unbiasedness, identification or finding a stationary point
\cite[pp.~316--317]{Fisher1922Foundations}.

\begin{hbdefinition}{def:p15-fe-properties}{Identification and a finite sufficiency criterion}
A target \(\tau(\theta)\) is identified by an experiment when
\(P_\theta=P_{\theta'}\) implies
\(\tau(\theta)=\tau(\theta')\). In a finite sample space, a
statistic \(T\) is sufficient when, for every value \(t\) having
positive probability, the conditional masses
\(P_\theta(Z=z\mid T=t)\) can be given by one kernel independent
of \(\theta\). Values impossible under a particular law do not
produce a conditional probability by dividing zero by zero.
This is the finite version of Fisher's conditional-information criterion,
not an assertion that every maximizing estimator satisfies it.
\end{hbdefinition}

Fisher compares the population-calibrated absolute-deviation and
root-mean-square calculations for a normal scale, then stresses that
population calibration need not mean equal precision. He also discusses
grouping corrections and the method of moments. His moment argument
requires the moments used to fit the curve to exist; variance control
for a fourth-moment calculation requires additional moment assumptions.
Calling a heavy-tailed distribution anomalous does not repair a moment
method that is undefined for it \cite[pp.~316, 318--322]{Fisher1922Foundations}.
The normal-scale comparison, with its limit variances, the efficiency
\(1/(\pi-2)\) of the mean error and Fisher's sufficiency-to-efficiency
argument, is \cref{thm:p15-w02b-p15-057-normal-scale,thm:p15-w02b-p15-057-suff-eff}.

\section{The actual binomial likelihood and its endpoints}

\begin{hbdefinition}{def:p15-fe-binomial}{Fisher's success/failure example}
Fix \(n\geq1\). The ordered observations \(Z_1,\ldots,Z_n\)
are independent Bernoulli variables with success probability
\(p\in[0,1]\). Put \(K=\sum_iZ_i\); the observed success and
failure counts are \(x=K\), \(y=n-K\).
The count mass is
\[
  P_p(K=x)=\binom nx p^x(1-p)^y.
\]
Removing only its parameter-independent binomial factor gives the
likelihood \(\mathcal L(p;x,y)=p^x(1-p)^y\).
The exponent-zero factors are one, including at an endpoint.
On \(0<p<1\), retain the paper's \(L\) for log likelihood:
\[
  L(p;x,y)=x\log p+y\log(1-p).
\]
The source's \(L=S(\log f)\) uses \(S\) as a summation sign,
not as our count statistic \(K\). The observed \(x,y\) are held
fixed during maximization \cite[pp.~323--324,331]{Fisher1922Foundations}.
\end{hbdefinition}

\begin{hbtheorem}{thm:p15-fe-binomial}{proved=annex:p15.fe.binomial}{Maximum, sampling precision and sufficiency}
On the closed interval the unique maximizer is \(\widehat p=x/n\).
When \(x,y>0\) it is interior and
\[
 L'(p)=\frac{x}{p}-\frac{y}{1-p},\qquad
 L''(p)=-\frac{x}{p^2}-\frac{y}{(1-p)^2}<0.
\]
On the open parameter interval there is no maximizer when \(x=0\)
or \(y=0\), only an endpoint supremum.
Under the stated independent model,
\[
 E_p\widehat p=p,\qquad
 \operatorname{Var}_p(\widehat p)=\frac{p(1-p)}n.
\]
For \(\epsilon>0\), its error probability is at most
\(p(1-p)/(n\epsilon^2)\), hence it is sampling-consistent.
The count, and therefore \(\widehat p\) at fixed \(n\), is sufficient.
For every possible count, the conditional law is uniform on the
\(\binom nx\) ordered strings having that count.
\end{hbtheorem}

This supplies an exact success case for Fisher's method. The variance
formula is not merely justified by the paper's curvature approximation;
it follows directly from independent Bernoulli sampling. The corrected
second derivative uses \((1-p)^2\), not the malformed
\(1-p^2\) denominator in the held p.~330 display. At boundary observations
a score equation cannot replace the global endpoint comparison.

\section{Coordinate invariance is not a prior probability}

Fisher uses \(\sin\theta=2p-1\) to show that a likelihood
maximizer is the same parameter value in different coordinates, whereas
declaring both coordinates equally probable selects different priors.
His criticism concerns obtaining a prior from ignorance alone; it does
not invalidate conditioning a specified joint probability law
\cite[pp.~324--327]{Fisher1922Foundations}.

\begin{hbtheorem}{thm:p15-fe-coordinates}{proved=annex:p15.fe.coordinates}{The source transformation and two explicit priors}
Let \(p=(1+\sin\theta)/2\) on
\(\theta\in[-\pi/2,\pi/2]\). Maximizing the transported likelihood
gives
\[
              \sin\widehat\theta=2\widehat p-1
                                    =\frac{x-y}{n}.
\]
The factor \(2n\) printed in Fisher's p.~325 solved equation is
incorrect under his own \(n=x+y\) convention.

A uniform probability density in \(\theta\) transports to
\(1/(\pi\sqrt{p(1-p)})\) on \(0<p<1\), whereas a uniform
prior in \(p\) transports to density \(\cos\theta/2\).
For \(x=y=1\), their respective posterior densities in \(p\) are
\[
 \frac8\pi\sqrt{p(1-p)},\qquad 6p(1-p).
\]
Both are proper and distinct. For \(x=1,y=2\), the uniform-\(p\)
posterior has its \(p\)-density mode at \(1/3\), but the same
probability measure's transported \(\theta\)-density mode corresponds
to \(p=3/8\). Its likelihood maximum still corresponds to \(p=1/3\).
\end{hbtheorem}

Likelihood does not become a parameter probability measure by drawing
its graph or assigning its maximum height one. Fisher's p.~327 language
that likelihood cannot be integrated must not be read as an assertion
that its mathematical function has no integral. For this example
\(\int_0^1p(1-p)\,dp=1/6\). The integral uses the chosen measure
\(dp\), which is not supplied by likelihood alone. Normalizing after
choosing that measure is precisely the additional prior choice, not a
coordinate-free deduction from the data.

Confidence instead concerns coverage of a random reported set under
each fixed parameter, as in \cref{def:p15-confidence-procedure}.
Posterior content concerns the specified joint law and prior. Curvature
or a relative likelihood height supplies neither coverage nor posterior
content without the corresponding sampling or probability construction.
Neyman's later procedure is a distinct presentation, not a theorem
implicitly present in Fisher's 1922 terminology.
\Cref{thm:p15-w02b-p15-057-three-statements} computes a curvature
interval, its exact fixed-parameter coverage and its posterior content in
one binomial model.

\section{What regularity can justify without an asymptotic shortcut}

The paper's curvature argument assumes a near-normal sampling law,
root-\(n\) localization, concentration of summed second derivatives,
and negligible higher Taylor terms. It initially disregards how the
putative standard deviation varies with the parameter. Those assumptions
cannot be consequences of writing \(L'=0\). A maximum might not exist,
might lie on a boundary, might not be unique, or might fail to concentrate
\cite[pp.~328--332]{Fisher1922Foundations}.

\begin{hbdefinition}{def:p15-fe-finite-regularity}{A finite common-support family}
Let \(\mathsf Z\) be finite, \(\eta\) range over an open real
interval, and \(q_\eta(z)>0\) be differentiable masses summing to one.
Define
\[
 U_\eta(z)=\partial_\eta\log q_\eta(z),\qquad
 J(\eta)=E_\eta U_\eta^2.
\]
For a fixed statistic \(T:\mathsf Z\to\mathbb R\), put
\(m(\eta)=E_\eta T\). Differentiation involves only finite sums.
For the curvature identity also assume the masses are twice differentiable.
These premises are stated here, not imported by calling a model regular.
\end{hbdefinition}

\begin{hbtheorem}{thm:p15-fe-score}{proved=annex:p15.fe.score}{Finite score identities and the precision bound}
Under \cref{def:p15-fe-finite-regularity},
\[
 E_\eta U_\eta=0,\qquad
 m'(\eta)=E_\eta[(T-m(\eta))U_\eta],\qquad
 \operatorname{Var}_\eta(T)J(\eta)\geq(m'(\eta))^2.
\]
If \(J=0\), then \(m'=0\); no division by zero is permitted.
For twice differentiable masses,
\(J=-E_\eta\partial_\eta^2\log q_\eta\).
Independent repetitions have information \(nJ\).
For any statistic of the \(n\) independent repetitions unbiased for
\(\eta\) throughout the open interval,
\(J>0\) gives variance at least \(1/(nJ)\).
The binomial proportion attains that bound at interior \(p\), with
one-trial information \(1/(p(1-p))\).
Its expected negative log-likelihood curvature is \(n/(p(1-p))\).
For interior observed \(\widehat p\), the negative curvature at that
maximum is \(n/(\widehat p(1-\widehat p))\).
\end{hbtheorem}

Expected information and observed curvature are different objects.
The first is an expectation under the fixed true sampling law; the
second is a function of recorded data. Replacing unknown \(p\) in
an exact variance by \(\widehat p\) gives a plug-in estimate, not
an identity for the true sampling variance or a confidence guarantee.

This finite precision theorem is an editorial result, not a claim that
Fisher printed a modern general information inequality. In continuous
families, replacing these finite sums by integrals needs hypotheses
justifying each differentiation: a common dominating measure/support,
appropriate derivative domination and finite score moments are examples
of relevant conditions. Parameter-dependent endpoints contribute boundary
terms. An asymptotic information formula additionally needs existence,
identification, localization and remainder/limit control. The explicit
local-and-tail Laplace conditions of \cref{def:p15-laplace-family}
illustrate why even an integration approximation requires more than a
stationary-point label; they are not themselves an MLE sampling theorem.
Hypotheses of exactly this kind are stated and used in
\cref{thm:p15-w02b-p15-057-information,thm:p15-w02b-p15-057-asymptotic}.

\section{The sufficiency argument and an exact countermodel}

Fisher partitions sample space into isostatistical regions where the
estimated parameters agree. In section~7 he observes that all sample
elements in such a region maximize their likelihood at the same value.
A factorization into a parameter/statistic factor and a conditional
parameter-free factor would indeed establish sufficiency. But agreement
of maximizing points does not establish that factorization. Fisher
acknowledges lack of a rigorously complete general proof on p.~323;
the summary's stronger description on p.~367 does not repair the step
\cite[pp.~323,330--331,367]{Fisher1922Foundations}.

\begin{hbtheorem}{thm:p15-fe-failures}{proved=annex:p15.fe.failures}{Maximization, identification and independence can fail separately}
For one observation \(Z\in\{1,2,3\}\), take \(\eta\in[0,1]\) and
\[
 q_\eta(1)=\frac{2-\eta^2}{5},\quad
 q_\eta(2)=\frac{2-\eta^2/2}{5},\quad
 q_\eta(3)=\frac{1+3\eta^2/2}{5}.
\]
These positive masses sum to one and identify \(\eta\).
The unique MLE statistic is \(T(1)=T(2)=0,T(3)=1\), but it is not
sufficient: the conditional probability of \(Z=1\) given \(T=0\)
is \(1/2\) at \(\eta=0\) and \(2/5\) at \(\eta=1\).

Extending the same family to \([-1,1]\) makes the signed target
unidentified. For independent repetitions of that family, no estimator
sequence can be consistent for every signed
parameter, because the laws at \(a\) and \(-a\) agree.

In a different experiment let every binary observation equal the same
Bernoulli variable \(B\) with \(0<p<1\). The proportion equals \(B\)
for every \(n\), has variance \(p(1-p)\), and is not consistent for
\(p\). Fitting the independent-binomial likelihood does not restore
the missing sampling premise or justify variance \(p(1-p)/n\).
These are constructed countermodels, not measured datasets.
\end{hbtheorem}

The first countermodel pinpoints the source's logical gap without relying
on densities at exceptional null observations. The second concerns what
the laws can identify, not optimization skill. The third concerns a
misspecified joint sampling model even though each marginal proportion
has the desired success probability. They diagnose different failures;
none is resolved by treating likelihood as a posterior.

\section{Fisher's heavy-tailed location example}

Fisher's random-line intersection example gives the unit-scale Cauchy
location density \(1/(\pi(1+(u-m)^2))\). The location \(m\) is
well defined even though a first-moment method is unavailable. He uses it
to show that a distribution's infinite variance need not make location
estimation impossible \cite[pp.~321--323]{Fisher1922Foundations}.

\begin{hbtheorem}{thm:p15-fe-cauchy}{proved=annex:p15.fe.cauchy}{The exact mean failure and a consistent median}
For independent observations from that Cauchy location law, the sample
average has the same law as one observation. Thus for \(\epsilon>0\),
\[
 P_m(|\overline X_n-m|\leq\epsilon)
                       =\frac2\pi\arctan\epsilon,
\]
independently of \(n\). The first absolute moment and the second
moment are infinite. Symmetric cancellation is not a finite expectation.

For odd \(n=2r+1\), let \(M_n\) be the middle order statistic and
\(\delta_\epsilon=\arctan(\epsilon)/\pi\). Then
\[
 P_m(|M_n-m|>\epsilon)
       \leq\min\!\left(1,\frac1{2n\delta_\epsilon^2}\right).
\]
Hence the median is sampling-consistent. This finite bound does not
claim the paper's more precise asymptotic median error or MLE efficiency.
The average-law proof uses the already supplied Cauchy characteristic
function result \cref{thm:p15-characteristic-inverse-examples}.
\end{hbtheorem}

Rejecting an extreme observation merely because it is extreme changes
the observation procedure. Fisher recommends considering whether the
error distribution is nonnormal instead of silently discarding observations.
That is a model/recording warning, not an instruction to ignore independently
documented recording errors or to adopt a Cauchy model for every dataset.

\begin{hbmethod}{def:p15-fe-use}{Applying a likelihood argument at its actual scope}
Specify the observation law, dependence structure, parameter space and
estimand. Check whether different parameter values give the same law.
Form the full likelihood and justify every discarded factor. Find a
global maximum or state that only a supremum exists; do not treat a
score equation as existence or uniqueness.

For information calculations state common support, derivative/expectation
conditions, moment finiteness and the statistic's parameter independence.
For sufficiency inspect the conditional law, not just the maximizer.
For large-sample claims supply localization, Taylor remainder and limit
control separately. For a confidence set or posterior, use its own
sampling-coverage or joint-law definition rather than likelihood height.
Keep the exact finite conclusions distinct from a source's asymptotic
approximation and from the adequacy of the chosen model.
\end{hbmethod}

\index{Fisher, R. A.!estimation criteria}
\index{likelihood!not a parameter probability}
\index{sufficiency!maximum likelihood counterexample}
\index{regularity!score and information}
